\hypertarget{appendix-8f6a01f284a7}{}
\section*{Conditioning and singularity}
\section{Conditioning and singularity}\label{conditioning-and-singularity}

\textbf{Curriculum depth:} Developed · \textbf{First seen:} C004 · \textbf{Latest development:} Phase7

\subsection{Definition and mechanism}\label{definition-and-mechanism}

Conditioning measures sensitivity to input perturbation. \(A=10^{-4}I_2\) has determinant \(10^{-8}\) but condition number one in the 2-norm. Uniform scaling changes determinant magnitude without changing this conditioning measure.

Phase 7 develops \(\kappa(A)=\|A\|\|A^{-1}\|\), near-parallel-column sensitivity, pivoted elimination and inverse residuals. Equal determinants can conceal very different stretch ratios: \(\operatorname{diag}(10^{-10},10^{-10})\) and \(\operatorname{diag}(1,10^{-20})\) both have determinant \(10^{-20}\) but conditions one and \(10^{20}\). \href{Continuous\%20Notes/Computing\%20Foundations\%20-\%20From\%20Bits\%20to\%20C++\%20and\%20Systems\%20APIs\#An\%20inverse\%20can\%20exist\%20and\%20still\%20amplify\%20uncertainty}{Full worked sensitivity example}.

Subtraction of nearby inputs exposes input uncertainty relative to a small difference. Sensitivity belongs to the problem; exact subtraction of stored operands cannot recover information lost when those inputs were rounded.

\subsection{Prerequisites}\label{prerequisites}

\begin{itemize}
\tightlist
\item
  \hyperref[appendix-0ec599f8ba27]{Numerical tolerance policies} --- An absolute numerical threshold can change with scale without measuring mathematical sensitivity.
\end{itemize}

\subsection{Used by}\label{used-by}

\begin{itemize}
\item
  \hyperref[appendix-48d2e66378ab]{Matrix inversion pivoting and residuals}
\item
  \hyperref[appendix-a7fbaf6d1b1c]{Forward and backward error} --- Input sensitivity is distinct from errors an implementation introduces.
\item
  \hyperref[appendix-a2773a6a2f9d]{Cancellation and loss of significance} --- Input sensitivity is distinct from errors an implementation introduces.
\end{itemize}

Connect later dependent concepts here as their lessons arrive. - \href{}{Gaussian elimination and pivoting (02-Foundations-II, appendix-f6e77f4a200f)} --- Extends this foundation into linear-systems and decomposition analysis. - \href{}{QR factorization and least squares (02-Foundations-II, appendix-c3421e6b4f2c)} --- Extends this foundation into linear-systems and decomposition analysis. - \href{}{ODE integration and stability (04-Visual-Computing-II, appendix-18eacf61a04b)}

\begin{itemize}
\tightlist
\item
  \href{}{A posteriori error bounds for contractive iteration (02-Foundations-II, appendix-d83b0ad1be7f)}
\item
  \href{}{Algorithmic stability (02-Foundations-II, appendix-aa2bf0d79639)}
\item
  \href{}{Bracketed and open root solvers (02-Foundations-II, appendix-f63eb850796d)}
\item
  \hyperref[appendix-31b8c295e05f]{Numerical error budgets}
\item
  \href{}{Stable quadratic root evaluation (02-Foundations-II, appendix-b9afba7f63e5)} \#\# Recall and next depth
\end{itemize}

Reconstruct the concrete example in the definition and identify its validity assumptions.

\textbf{Pending:} Robust condition estimation, growth factors, advanced factorizations and application-specific error budgets extend the current inversion laboratory.

\subsection{Source and navigation}\label{source-and-navigation}

\href{Continuous\%20Notes/Computing\%20Foundations\%20-\%20From\%20Bits\%20to\%20C++\%20and\%20Systems\%20APIs\#When\%20arithmetic\%20meets\%20uncertainty}{M005 development}

\href{Continuous\%20Notes/Computing\%20Foundations\%20-\%20From\%20Bits\%20to\%20C++\%20and\%20Systems\%20APIs\#A\%20stretching\%20ruler\%20for\%20real-valued\%20quantities}{M004 teaching chapter} · \hyperref[appendix-fd090d982a15]{Per-part definitions} · \href{}{Dependency Map (02-Foundations-II, appendix-4ad47bdd269a)}
